% Generated draft: requires mathematical and source review.
\documentclass[11pt]{article}
\usepackage[margin=25mm]{geometry}
\usepackage{fontspec}
\setmainfont{Latin Modern Roman}
\usepackage{amsmath,amssymb,mathtools,unicode-math}
\setmathfont{Latin Modern Math}
\usepackage{enumitem,needspace,xurl,hyperref}
\setlength{\parindent}{0pt}
\setlength{\parskip}{5pt}
\setlength{\emergencystretch}{3em}
\newcommand{\N}{\mathbb{N}}
\newcommand{\Z}{\mathbb{Z}}
\newcommand{\Q}{\mathbb{Q}}
\newcommand{\R}{\mathbb{R}}
\newcommand{\C}{\mathbb{C}}
\newcommand{\F}{\mathbb{F}}
\newcommand{\A}{\mathbb{A}}
\newcommand{\PP}{\mathbb{P}}
\newcommand{\E}{\mathbb{E}}
\newcommand{\eps}{\varepsilon}
\newcommand{\dd}{\,\mathrm{d}}
\begin{document}
\section*{UnsolvedMath ID: 30004893}
\textbf{Source problem code:} catalogue entry
\subsection{Definitions and mathematical statement}
The supplied record concerns Conjecture 4 in the cited Oberwolfach report. Its intended category \(\mathcal C\) consists of finite-dimensional representations of a quantum affine algebra. The record does not specify the algebra, its deformation parameter, the admissible categories, or the hypotheses needed for the constructions below. Consequently, the following is a conditional formulation of the supplied target, requiring the source definitions before it becomes a fully specified conjecture.

For an exact category \(\mathcal C\), its Grothendieck group \(K(\mathcal C)\) is generated by isomorphism classes \([V]\), subject to \([V]=[U]+[W]\) for each short exact sequence \(0\to U\to V\to W\to0\) in its specified exact structure. The required exact and tensor structures on the intended \(\mathcal C\) are not supplied. Denote the source's parameter-dependent Grothendieck module by \(M=K(\mathcal C)_t\), its coefficient ring by \(R\), and its quantum Grothendieck ring by \(Q=K_t(\mathcal C)\). These are distinct source-designated objects; the record does not define their coefficient ring, parameter conventions, or underlying modules. Let \(*:M\times M\to M\) denote the bilinear operation that the source constructs from a monoidal Jantzen filtration. The filtration and the formula for this operation are absent from the record.

\textbf{Conditional conjectural target.} For every category covered by the intended source hypotheses, the operation \(*\) is associative:
\[
(x*y)*z=x*(y*z)\qquad(x,y,z\in M).
\]
Together with the intended anti-involution \(\iota\), it makes \(M\) an algebra with anti-involution. In particular, \(\iota\) is additive and satisfies
\[
\iota^2=\operatorname{id}_M,
\qquad
\iota(x*y)=\iota(y)*\iota(x)
\qquad(x,y\in M).
\]
The record does not specify whether \(\iota\) is \(R\)-linear or semilinear for an involution of \(R\), nor whether unitality is part of the asserted algebra structure.

There is an isomorphism of algebras with anti-involution
\[
\Phi:(K(\mathcal C)_t,*,\iota)\xrightarrow{\sim}(K_t(\mathcal C),\cdot,\iota_Q)
\]
which identifies the corresponding standard bases. Thus, under the source's intended coefficient conventions,
\[
\Phi(x*y)=\Phi(x)\Phi(y),
\qquad
\Phi\circ\iota=\iota_Q\circ\Phi.
\]
For each prescribed pair of corresponding standard bases, written schematically as \((b_\lambda)_{\lambda\in\Lambda}\) in \(M\) and \((s_\lambda)_{\lambda\in\Lambda}\) in \(Q\), the required identification is \(\Phi(b_\lambda)=s_\lambda\) for every \(\lambda\in\Lambda\). The bases, their indexing, their correspondence, and \(\iota_Q\) are not specified in the supplied record; this notation supplies no additional choices.

The supplied dated curation reports the conjecture proved for simply-laced quantum loop algebras and additional special categories, while reporting associativity and the claimed isomorphism conjectural in the general non-simply-laced setting. These are supplied status claims, not independently verified mathematical conclusions.
\subsection{Short English statement}
The conjecture asks whether the product obtained from a monoidal Jantzen filtration is associative and yields the quantum Grothendieck ring, respecting anti-involutions and the corresponding standard bases. The supplied record omits the precise category hypotheses and construction conventions, so this is a conditional statement requiring source consultation and human mathematical review. Its partially solved status is supplied dated curation.
\subsection{Sources}
\begin{itemize}
\item Supplied local UnsolvedMath registry record, permanent ID \(30004893\), source problem code \texttt{catalogue entry}, title \emph{Quantum Grothendieck Rings Beyond Simply-Laced Types}, in the Oberwolfach Reports Open Problems collection. The proposed year is recorded as 2021; no individual proposer is supplied.
\item Supplied report citation: \emph{Enveloping Algebras and Geometric Representation Theory} (2022), Oberwolfach Reports, pp.\ 2827--2891, DOI \texttt{10.4171/owr/2021/52}. \url{https://doi.org/10.4171/owr/2021/52}. The original-statement field refers to Conjecture 4 but does not reproduce its mathematical content.
\item Supplied literature reference: Fujita--Hernandez, \emph{Advances in Mathematics}. \url{https://doi.org/10.1016/j.aim.2026.110963}.
\item Supplied literature reference: Fujita--Hernandez, \emph{Monoidal Jantzen filtrations}. \url{https://arxiv.org/abs/2402.13544}.
\item Literature assessment dated exactly \texttt{2026-08-22}: the supplied curation labels the problem partially solved, reports proofs for simply-laced quantum loop algebras and additional special categories, and reports the general non-simply-laced associativity and isomorphism claims as conjectural. This assessment is supplied dated curation, not independently verified; its time-sensitive status and cited claims require independent verification.
\item The registry labels the statement \texttt{corrected\_verified} and states that the official report was checked directly and the clean statement restored definitions while removing extraction artifacts. This is supplied verification metadata; no independent source visit or verification is asserted here.
\end{itemize}
\end{document}
